\documentclass{article}
\usepackage[margin=1in]{geometry}
\usepackage[style=apa,backend=biber]{biblatex}
\addbibresource{/home/nicoluarte/Documents/repos/phd_DCCS/references/refs.bib}
\usepackage{titlesec}

% Tighten section spacing to save space elsewhere
\titlespacing*{\section}{0pt}{1.5ex plus 1ex minus .2ex}{1ex plus .2ex}
\titlespacing*{\subsection}{0pt}{1.5ex plus 1ex minus .2ex}{1ex plus .2ex}

\title{A category theory informed specification of the cerebellar computational model: causal test of the universal cerebellar transform}
\author{Luis Luarte}
\begin{document}
	\maketitle
	\begin{abstract}
		Cerebellar research is polarized between the \textbf{Universal Cerebellar Transform (UCT)} \cite{itoMovementThoughtIdentical1993}, which posits a uniform algorithm, and \textbf{multimodal parcellation} \cite{diedrichsenUniversalTransformMultiple2019}, implying functional heterogeneity. Current models fail to reconcile these, often succumbing to ad-hoc parameter tuning or the ``projector trap'' due to lacking formal specification.
		
		This project resolves this using \textbf{Category Theory (CT)}. Unlike standard modeling, CT imposes strict axioms (Identity, Composition) that eliminate biologically implausible architectures. We employ an ``axiomatic pass'': forcing the biological circuit to satisfy these axioms mathematically necessitates a predictive Forward Model. This replaces curve-fitting with structural proof derived from logical necessity.
		
		To bridge this specification with data, we define a \textit{functor} mapping internal categorical states to EEG metrics. We target Theta power as a proxy for error-signal processing (Climbing Fiber) and Gamma power for state-prediction binding (Purkinje). Finally, we execute a causal test applying Transcranial Magnetic Stimulation (TMS) to Crus I/II during two orthogonal tasks: motor force-tracking and an economic two-armed bandit. By tracking if TMS disruption propagates through Theta/Gamma dynamics invariantly across domains, we provide a causal proof of the UCT.
	\end{abstract}
	
	\section{Project proposal}
	\subsection{Justification}
	Our methodology is grounded in the \textbf{Category-Informed Model Specification (CIMS)}. This framework formally constructs the cerebellum as a neuro-category ($\mathcal{C}_{Cb}$), defining neuronal populations (Objects) and synaptic transfer functions (Morphisms). This moves beyond descriptive biology to establish a rigorous mathematical structure where composition rules are explicitly defined.
	
	To bridge this formalism and empirical observation, we define a \textbf{measurement functor}, $F: \mathcal{C}_{Cb} \to \mathcal{C}_{EEG}$. This functor translates hidden computational states into observable metrics, such as mapping internal "prediction error" to specific EEG oscillatory dynamics. To address scalp EEG noise, we employ \textbf{Bayesian model adequation}. We use the strict axioms of CIMS to define informative priors, allowing us to invert the functor and infer latent activity of specific morphisms (e.g., error-driven updates) directly from source-localized data.
	
	With this validated model, we define the motor force-tracking and economic decision-making tasks as \textbf{isomorphic input categories}. By equating the force vector (motor) with the reward probability vector (economic), we ensure both tasks engage the same internal morphisms. The critical test of the UCT follows: if the cerebellum utilizes a universal model, TMS perturbation must propagate through the functor identically in both tasks. We predict an \textbf{isomorphic collapse} of the commutative diagram, where the failure mode in the motor EEG is mathematically identical to that in the economic EEG. This distinguishes genuine universal computation from correlated deficits.
	
	\subsection{Literature review}
	
	\subsubsection*{1. The Universal Cerebellar Transform (UCT) and Cognitive Extension}
	The UCT hypothesis, originally formulated by Ito and expanded by Schmahmann, posits that the uniform cytoarchitecture of the cerebellar cortex implies a uniform computational operation performed on all inputs, regardless of their modality \parencite{itoMovementThoughtIdentical1993, schmahmannDysmetriaThoughtClinical1998}. While traditionally viewed as a motor error-correction machine \parencite{ecclesCerebellumNeuronalMachine1967}, recent decades have witnessed a paradigm shift recognizing the cerebellum's critical role in cognitive domains, including language, social cognition, and economic decision-making \parencite{koziolConsensusPaperCerebellums2014}.
	
	Specific to economic domains, evidence implicates the lateral cerebellum (Crus I/II) in \textbf{reinforcement learning} and \textbf{probabilistic reversal learning}. Recent human studies demonstrate that the cerebellum encodes \textit{Reward Prediction Errors} (RL-PEs), the discrepancy between expected and actual outcomes, which are critical for updating economic value functions \parencite{huvermannCerebellumContributesPrediction2025}. Crucially, while cerebellar damage often spares the initial acquisition of value associations, it specifically impairs \textbf{reversal learning}, the ability to flexibly update these associations when reward contingencies change \parencite{thomaCerebellumInvolvedRewardbased2008}.
	
	For instance, combined TMS and EEG studies have shown that disrupting the lateral cerebellum blunts the \textbf{Feedback-Related Negativity (FRN)}, a cortical proxy for prediction error processing, effectively decoupling the cerebellum from the cortical economic monitoring network \parencite{huvermannCerebellumContributesPrediction2025}. This literature provides the empirical justification for selecting a \textbf{Two-Armed Bandit} task as the cognitive isomorphism to the standard visuomotor force-tracking task: both tasks do not merely require memory, but the sequential, error-driven estimation of a hidden state variable (force vs. reward probability) regulated by the same computational topology \parencite{schultzNeuronalRewardDecision2015}.
	
	\subsubsection*{2. The Challenge of Multimodal Parcellation}
	Despite the UCT's theoretical elegance, it faces significant challenges from \textbf{Multimodal Parcellation} models. Recent functional imaging studies argue that the cerebellum is functionally heterogeneous, divided into distinct "parcels" that mirror their specific cortical inputs \parencite{diedrichsenUniversalTransformMultiple2019}. This view suggests that "motor" cerebellum and "cognitive" cerebellum are distinct processors, supported by evidence that separate cerebellar sub-regions activate during narrative language processing versus motor action.
	
	The conflict between these two models, Unity vs. Parcellation, remains unresolved because current computational models lack formal rigor. Standard models often treat the Granule Layer as a generic ``random projection'' engine or Reservoir Computer \parencite{souresSpikingReservoirNetworks2019}, a description that is mathematically loose enough to fit either hypothesis depending on parameter tuning. This ambiguity leads to the ``projector trap'', where the functional specificity of the cortical input is mistaken for specificity in the cerebellar computation itself.
	
	\subsubsection*{3. Category Theory in Neuroscience: From Algebra to Axioms}
	To resolve this specification gap, the field is increasingly turning to \textbf{Category Theory (CT)}. At its core, CT is the mathematical study of structure, defined not by what components \textit{are}, but by how they \textit{compose}. To formalize the cerebellum, we utilize three fundamental concepts:
	\begin{enumerate}
		\item \textbf{Objects:} The state-spaces of neuronal populations (e.g., Granule Cell vectors).
		\item \textbf{Morphisms:} The structural transformations (synaptic weights) between objects.
		\item \textbf{Monads:} A structure that encapsulates ``computational context.''
	\end{enumerate}
	
	Recent frameworks in ``Compositional Neuroscience'' have modeled the cerebellum as a \textbf{Context Monad} (or Reader Monad). Intuitively, a Monad is a wrapper that adds capabilities to a simple value. The Granule Layer acts as a Context Monad because it does not just relay the Mossy Fiber input ($x$); it ``lifts'' it into a richer, high-dimensional state space ($T(x)$) where the input is bound to a specific temporal or multimodal context \parencite{ghoshCompositionalArchitectureImplicit2025}.
	
	However, there is a critical limitation: Monads describe the \textit{architecture} (the static wiring), but not the \textit{dynamics} (the learning). Knowing that the Granule Layer is a Monad tells us \textit{where} the information goes, but not \textit{how} the error signal corrects it. Our proposal improves on this by applying an \textbf{Axiomatic Pass}. We do not just assume a monadic structure; we prove that for the system to maintain the axiom of \textbf{Compositionality} ($h = g \circ f$) under perturbation, it is mathematically necessitated to implement a predictive Forward Model. This moves the field from \textit{descriptive algebra} (labeling the circuit) to \textit{mechanistic necessity} (proving the learning rule is a requirement of the structure) \parencite{tamimCategoricalInvariantsLearning2025}.
	
	\subsubsection*{4. Validation via Bayesian Model Adequation and Functors}
	A critical limitation in computational neuroscience is validating these abstract models against noisy biological data. Current validation efforts typically employ \textbf{Representational Similarity Analysis (RSA)} \parencite{kriegeskorteRepresentationalSimilarityAnalysis2008}. While effective for static representations, RSA suffers from the ``mimic effect'': mechanistically distinct systems can yield highly correlated representational geometries, rendering RSA insufficient for validating specific algorithms \parencite{dujmovicObstaclesInferringMechanistic2022}. Furthermore, RSA lacks the granularity to validate the \textit{process} of computation itself \parencite{kriegeskorteInferringBraincomputationalMechanisms2016}.
	
	To address this, we propose integrating Category Theory with \textbf{Bayesian Model Adequation}. By defining a \textbf{Measurement Functor} $F: \mathcal{C}_{Cb} \to \mathcal{C}_{EEG}$, we treat the mapping from neural states to electrophysiological signals as a structure-preserving transformation. This allows us to invert the functor using Bayesian priors derived from the categorical axioms. For example, if the category specifies that an error signal must update a weight morphism, the Bayesian model can infer the probability of this specific computational event given the observed Theta-band power, filtering out noise that does not conform to the axiomatic structure \parencite{gavranovicLearningFunctorsUsing2020}.
	
	\subsubsection*{5. Empirical Measurement: EEG and TMS in Crus I/II}
	The operationalization of this framework relies on recording cerebellar activity, historically a challenge due to the closed-field geometry. However, recent literature confirms that high-density EEG can detect cerebellar signals using source localization constrained to the posterior fossa. Specifically, ``Topological Reinforcement Learning'' algorithms (like TOREADA) have demonstrated the feasibility of extracting decision-related manifolds from EEG data \parencite{kalitzinTopologicalReinforcementAdaptive2024}.
	
	Furthermore, TMS studies have successfully established causal links between the lateral cerebellum (Crus I/II) and reinforcement learning functions \parencite{huvermannCerebellarSinglePulseTMS2025}. Our protocol leverages these findings, using MRI-guided TMS to target specific ``cognitive'' nodes, while the functor ensures that the resulting electrophysiological perturbations are interpreted through a formal computational lens rather than simple signal attenuation.
	
	\subsection{Objectives}
	
	\subsubsection*{General objective}
	To develop and validate a \textbf{Category-Informed Model Specification (CIMS)} of the cerebellar microcircuit, and to utilize this framework to demonstrate—via causal TMS perturbation—that the computational algorithm remains invariant across orthogonal motor and economic domains, validating the UCT.
	
	\subsubsection*{Specific objectives}
	\begin{enumerate}
		\item \textbf{Formal Construction of $\mathcal{C}_{Cb}$:} Mathematically define the cerebellar circuit as a category (Objects: populations; Morphisms: synapses) and apply an ``Axiomatic Pass'' to demonstrate that satisfying Identity and Composition axioms necessitates the emergence of a predictive Forward Model, establishing the structural ``Ground Truth''.
		
		\item \textbf{Operationalization of Measurement Functor ($F$):} Construct and validate a functor mapping internal states to EEG metrics (Prediction Error $\to$ Theta; State Prediction $\to$ Gamma) using Bayesian Model Adequation to ensure observed spectral changes represent specific computational morphisms.
		
		\item \textbf{Causal Interrogation via Orthogonal Tasks:} Apply continuous Theta Burst Stimulation (cTBS) to Crus I/II during mathematically isomorphic Visuomotor Force-Tracking and Two-Armed Bandit tasks to isolate the intrinsic circuit transform from input modality.
		
		\item \textbf{Verification of Isomorphic Diagram Collapse:} Analyze TMS-induced signal degradation using Equivalence Testing to demonstrate that the breakdown of the commutative diagram is mathematically identical across tasks, thereby proving the UCT.
	\end{enumerate}
	
	\subsection{Methods}
	
	\textbf{Experimental Design and Participants:} We will employ a within-subject, double-blind, sham-controlled crossover design with a sample of $N=30$ healthy adults (18-40 years). This sample size is calculated to achieve 80\% power ($\alpha=0.05$) to detect a medium effect size ($d=0.5$) in paired comparisons, accounting for potential attrition. The experiment follows a $2 \times 2$ factorial arrangement: Task Domain (Motor vs. Economic) $\times$ Stimulation (Active cTBS vs. Sham). All participants will complete both tasks under both stimulation conditions across four separate sessions, spaced by a washout period of at least seven days to prevent carry-over effects of the neurostimulation.
	
	\textbf{Variables and Procedures:} The Independent Variable is the application of continuous TMS Theta Burst Stimulation (cTBS) targeted via T1-weighted MRI neuronavigation to the left lateral cerebellum (Crus I/II). The Dependent Variables are strictly operationalized through our Category-Informed Model Specification (CIMS) to ensure cross-domain comparability. Behavioral metrics include \textit{Error Magnitude} (Euclidean distance from target) for the Motor task and \textit{Regret} (probability of suboptimal choice) for the Economic task. Simultaneously, high-density EEG (64-channel) will capture the Neural Dependent Variables defined by our Measurement Functor: Theta-band power (4--8 Hz) as a proxy for prediction error signaling (climbing fibers) and Gamma-band synchrony (30--80 Hz) representing state prediction maintenance (Purkinje cells).
	
	\textbf{Statistical Analysis:} Analysis will focus on testing the hypothesis of Isomorphic Collapse. We will first preprocess EEG data using LORETA source localization constrained to the cerebellar cortex. We will then calculate the "perturbation vector" ($\Delta$)—the difference between Active and Sham conditions—for both neural and behavioral metrics. Using Two One-Sided Tests (TOST) for statistical equivalence, we will test whether the perturbation vector in the Economic domain is statistically equivalent to the perturbation vector in the Motor domain (after domain-specific normalization). A rejection of the null hypothesis of difference ($p < 0.05$) in favor of equivalence will confirm that the TMS-induced "virtual lesion" degraded a single, domain-general computational operator, validating the Universal Cerebellar Transform.
	
	\subsection{Expected results}
	
	\subsubsection*{1. Formal Derivation of the Computational Model}
	We anticipate the neuro-category $\mathcal{C}_{Cb}$ will pass axiomatic verification, yielding a \textbf{Cerebellar Homomorphism Theorem}. \textbf{Mathematically}, by enforcing \textit{Compositionality} on the morphism chain, we expect to show that maintaining \textit{Identity} under perturbation requires the internal synaptic morphism to approximate the inverse of external dynamics. \textbf{Resulting Model:} This will define the Deep Nuclei as implementing a strict \textbf{Comparator Functor} ($Output = Drive - \hat{Prediction}$), possessing fewer free parameters than reservoir models as connectivity is fixed by categorical structure.
	
	\subsubsection*{2. Bayesian Adequation of the Measurement Functor}
	We expect Bayesian Model Adequation to favor the CIMS over null models. \textbf{Posterior Check:} Observed Theta (4--8 Hz) and Gamma (30--80 Hz) densities should fall within the 95\% Credible Interval of the posterior predictive distribution of $\mathcal{C}_{Cb}$. \textbf{Model Comparison:} We anticipate the CIMS will show lower WAIC scores than generic GLMs, specifically capturing the \textit{temporal variance} of Theta-Gamma coupling, validating the physical instantiation of the "Prediction Error" morphism.
	
	\subsubsection*{3. TMS Perturbation Metrics and Isomorphism}
	We define precise metrics for the unperturbed and perturbed states. \textbf{Baseline:} We expect event-locked Theta bursts following errors and sustained Gamma during state maintenance in both tasks. \textbf{Perturbation:} cTBS is hypothesized to induce a \textbf{Frequency-Specific Dampening} ($p < 0.01$) of the error-related Theta burst in \textit{both} tasks. While standard spatial RSA may show low similarity, our Dynamic Causal Modeling is expected to show high \textbf{Algorithmic Similarity}, proving the error integration \textit{process} is identical.
	
	\subsubsection*{4. Statistical Testing of Isomorphic Collapse}
	We will test the UCT via \textbf{Equivalence Testing (TOST)}. Defining $\vec{v}_{motor}$ and $\vec{v}_{econ}$ as the vectors of change (Post-TMS minus Pre-TMS), we predict rejecting the null hypothesis of difference ($H_0: |\vec{v}_{motor} - \phi(\vec{v}_{econ})| > \Delta$) in favor of equivalence. This \textbf{Isomorphic Collapse} would confirm that the functional degradation in the economic task is a strict mathematical transformation of that in the motor task, verifying the disruption of a domain-general operator.
	
	\subsection{Proposal novelty}
	This proposal pioneers the CIMS, shifting computational neuroscience from curve-fitting to axiomatic necessity. We perform the first ``Axiomatic Pass,'' proving a Forward Model is a structural requirement of biological axioms. We innovate by bridging formalism and data via a Bayesian Measurement Functor, translating noisy EEG into verifiable computational states. This enables the first causal, domain-invariant test of the UCT via Isomorphic Collapse.
	
	\section{Project timeline}
	The project spans 24 months (Jan 2026 -- Jan 2028).
	
	\subsection*{Year 1: Specification and Adequation (2026)}
	\textbf{Obj 1: Formal Construction (Months 1--5):} Jan--Mar: Math definition and ``Axiomatic Pass''. Apr--May: In-silico stability simulations.
	
	\noindent \textbf{Obj 2: Functor Operationalization (Months 6--11):} Jun--Jul: Task programming and Pilot EEG ($N=5$). Aug--Oct: Bayesian code dev and prior construction. Nov--Dec: Functor validation.
	
	\subsection*{Year 2: Causal Interrogation (2027)}
	\textbf{Obj 3: TMS Interrogation (Months 12--19):} Jan--Feb: Recruitment/MRI. Mar--Aug: cTBS Data Acquisition (Active vs. Sham).
	
	\noindent \textbf{Obj 4: Isomorphic Proof (Months 20--24):} Sep--Oct: LORETA/Functor Inversion. Nov--Dec: Statistical Equivalence Testing. Jan 2028: Manuscript prep.
	
	\newpage
	
	\printbibliography
	
\end{document}